% ch02 Introduction to finance (书页7-24 / p025-042)
\chapter{金融学导论}

本章是全书的“金融学地图”：它不推导任何具体的定价公式，而是把后文需要的最小编号集合一次配齐——股票价格$S(t)$与波动率$\sigma$（第3--8章的主角）、即期利率$r(t)$与远期利率$f(t,x)$、零息与附息国债价格$P(t,T)$、$\mathcal{B}(t,T)$（第9--15章的主角），以及贯穿全书的三个原理：有效市场假说（回答“模型为什么必须是随机的”）、货币的时间价值（回答“为什么要贴现”）、无套利与鞅测度（回答“为什么贴现期望就是价格”）。本章编号公式只有10个，但式(2.1)、(2.4)、(2.6)、(2.10)组成的“贴现期望定价公式族”是全书反复调用的引擎。阅读本章只需初等概率论；鞅的严格定义见原书附录A.1，完备市场中风险中性测度的存在与唯一性见原书附录A.6。

\section*{开篇：资本、金融资产与金融工具}

原书正文（书页7--8）先划定学科版图。经济学关心维系社会物质与精神生活的各种生产活动，其实物形态是实物资产（real assets）——资本品、管理、劳动力等；资本（capital）指一个社会实物资产的经济价值。在多数发达经济中，经济资产取货币化形态，资本可以被赋予货币价值即纸面形态，称为资本的货币形式（money form of capital）。金融学（finance）就是研究资本货币形式的经济学分支，不确定性与风险在其中具有根本重要性。本书聚焦金融资产（financial assets）：与实物资产相对，它们是“纸”，持有人凭其持有对一部分实物资产及其产生收入（若有）的索取权——例如股票持有人有权分得年度红利，并在公司清算时按比例分得剩余资产。金融资产又称证券（security），其具体形态（股票或债券等）称金融工具（financial instrument）；同一张纸对发行方而言同时是金融负债（financial liability），因为其利润与资产要在全体股东之间分割。

原书在此强调金融学区别于其他经济学分支的一点：它首先是实证学科（empirical discipline）。金融业每天产生海量数据，另有堆积如山的历史数据；市场强迫模型“实用、透明、可校准、可检验”。这使金融比其他经济学分支更接近自然科学——也正是本书敢于引入物理数学工具的经验论基础。

一个关键的供给结构：股票与债券处于正净供应（positive net supply）；衍生品则处于零净供应（zero net supply）——持有人数目恰等于卖出人数目，是纯粹的零和博弈，持有人的收益恒等于卖出方收益的相反数。这一点预先决定了衍生品的定价逻辑：不能靠“供求均衡”，只能靠无套利复制（见本章2.5、2.6节）。

原书把金融工具归为三种初级形态。其一，股权（equity），即普通股与股份：代表公司所有权份额，不保证回报，只有公司盈利时宣布红利的按比例份额；股价随公司经营起落，故股东暴露于公司风险之下，但只负有限的责任（limited liability）——最多损失初始投资，因此股票价值非负，最小值为零；股权是资产，因为持有人是资本的净所有者。其二，固定收益证券（fixed income securities），即债券（bond）：由公司与政府发行，承诺单笔或一列固定支付；债券是债务工具，发行人实为向买方借款，偿付安排在固定的时间区间上，该区间称为债券的到期日（maturity）。债券种类繁多：例如五年期附息美国国债承诺每半年付息、五年末还本，而零息美国国债只在到期日一次性付清。债券常被认为比股票安全——只要发行人不破产且持有人能持有到期，回报有保证；但由于利率风险、信用风险以及外币债券的汇率风险，债券组合的亏损完全可能不亚于甚至超过股票组合。其三，衍生证券（derivative securities）：收益派生自其他金融资产——例如取决于某股票或另一个衍生品的价格。

三种初级形态可近乎无穷地组合出新工具：例如优先股（preferred stock）兼有股权与债务特征——既分得发行人股权的份额，又像债券一样收取一列（固定）支付。原书随后给出理论金融的三大问题：货币的时间价值、风险与收益、证券与资产的估值；而为投资者量身创造新工具并做无套利定价，其设计、风险—收益分析与对冲构成金融工程（financial engineering）——这正是本书前半（期权与对冲）与后半（利率与固定收益）两条主线的总纲。

\section{有效市场：证券的随机演化}

本节要解决的问题是：股票价格凭什么应当用随机过程来描写？原书的回答不是“数据看起来像随机”，而是把随机性当作有效市场假说（efficient market hypothesis）的逻辑推论：一旦市场消化了全部可得信息并达到均衡，剩下的价格变动就只能是噪声。

先补齐市场词汇。买入（或约定买入）金融资产称持有多头（long position）或做多（going long）；卖出称做空（short position，shorting）；卖出自己并不实际持有的资产称卖空（short selling），回购日通常定在未来某时。市场组织有三种基本形态：直接市场（direct market，买卖双方直接搜寻成交）、经纪市场（brokered market，经纪人收费撮合）与拍卖市场（auction market，买卖双方在集中化市场中同时互动）。

有效市场（efficient market）概念与证券的“公平价格”（fair price）绑定：市场均衡时证券取得公平价格，投资者不再交易；均衡中的有效市场，其工具价格对公平价格只有小而暂时的偏离。有效市场又与市场信息紧密相关——区分不同市场参与者的正是各自可得的市场信息量。市场信息分三类：(a)金融资产价格与收益的历史数据，(b)关于金融资产各方面的公开数据，(c)少数参与者私有的信息；据此可分别定义市场效率的弱式（weak）、半强式（semi-strong）与强式（strong）形态。直观地说，有效市场意味着大多数买卖双方财富与信息对等，没有任何一方占（不公平的）优势。原书给出有效市场假说的精确表述，其内容可转述为：对处于均衡的金融市场，没有任何参与者在既有禀赋与信息下想再交易；价格揭示可得的市场信息；新信息以随机的、零碎的方式流入，由于信息本身不完备，市场参与者做出随机响应，于是各金融工具的价格随机变动。

这里有一个影响深远的推论：市场一旦均衡，所有证券的价格变化（除漂移（drift）外）都是随机的。理由是：有效市场中工具价格已经吸收了全部市场信息并形成均衡价格，此后任何对均衡的偏离都应当不确定、不可预测，且向上与向下变动的可能性相等。于是价格变化必须用随机变量（random variable）表示：设某股权在时刻$t$的价值为$S(t)$，其变化率$\mathrm{d}S/\mathrm{d}t$被建模为随机变量，这又蕴含$S(t)$本身也是随机变量，只在初始时刻$t_0$给定确定性初始条件；随机性的强度由证券的波动率（volatility）刻画，记作$\sigma_S$或简记$\sigma$。波动率越大，价格的随机涨落越大；$\sigma=0$则未来演化完全确定。\footnote{这一论证是启发式的。它的严格化就是后续章节的随机微分方程$\mathrm{d}S=\mu S\,\mathrm{d}t+\sigma S\,\mathrm{d}W$的形式：漂移$\mu$在2.3节以风险溢价的面目回归，$\sigma$从第3章起进入Black--Scholes与Merton--Garman方程。}所有投资者面对的风险，正是金融工具随机演化的反映，而后者最终反映（金融）市场吸纳底层实体经济全部相关特征的方式。

要紧的是，有效市场假说并不主张消息不能移动市场。它说的是：未预期的、出乎意料的新信息会扰乱各证券市场价格的均衡，并把它们系统地搬移到一组新均衡价格；一旦均衡达成，普通信息几乎人人可得，于是只造成所揭示价格的随机变动。

该假说可否被经验检验？原书引文献[23]指出，有效市场的存在暗中包含两个假说——效率假说本身，加上“市场处于某个特定均衡”的假说；能够被经验检验的只是二者的联合假说（joint hypothesis）。围绕这一联合检验，学术界关于市场效率的争论经久不息。

原书以一个物理类比收束本节（书页11首段）：市场均衡类似热力学系统的平衡——平衡态中单个粒子的位置与速度仍然随机，正如均衡中金融工具的价格仍然随机；市场效率则类似热机效率：没人指望真实热机达到$100\%$的效率，$60\%$--$70\%$已属常见；同样，即使金融市场并非完全有效，应用基于有效市场概念的数学模型往往仍然成立。

\section{金融市场}

本节回答两个问题：金融工具在哪里、以何种组织形式交易？市场的体量有多大？

金融市场的第一层划分是资本市场（capital markets）与货币市场（money markets）。资本市场交易初级形态的金融工具——股权、债务与衍生品；指数（indexes）本是一篮子证券的加权平均、属于资本市场，但因其重要性与深度被当作独立实体处理。货币市场严格说属于债务市场，但其工具高度流动、期限极短（现金及现金等价物、外汇交易等），故另设市场。原书给出如下分类（转述）：(1)资本市场——股权市场（普通股、优先股），债务市场（国债（政府）票据与债券、公司债与市政债、抵押贷款支持证券（mortgage-backed securities，MBS）），衍生品市场（期权、远期与期货）；(2)指数——股权指数（道琼斯与标普指数、日经指数、DAX指数、STI指数等），债务指数（债券市场指标）；(3)货币市场——现金定期存款、短期国库券（Treasury bills）、大额存单（certificates of deposit）、商业票据（commercial paper）、欧洲美元存款（Eurodollar，指存于非美银行或美国银行海外分行的美元存款）。

\paragraph{原书图2.1（1993年美国资本市场构成，总值\$11.3万亿）}图饼显示：股权仅占$36\%$，债务市场合计占$64\%$——国债$26\%$、MBS $14\%$、公司债$11\%$、市政债$11\%$、高收益债$2\%$。若再把货币市场计入，股权份额还要更低。

\paragraph{原书图2.2（1993年全球债务市场的借款人构成）}全球债务市场在正文中的数字为1993年\$14.08万亿，并称“图2.1[100]给出主要国际借款人”——实际上给出借款人构成的是图2.2，其图题写的是\$14.8万亿。正文与图题在交叉引用与数字上均不一致（原书排印疏漏，两说照录）。借款人构成：美国$46\%$、日本$18\%$、德国$10\%$、英国$8\%$、其他$8\%$、意大利$5\%$、法国$5\%$。

量级感（原书照录的数据）：美国GDP（2001年）约\$10万亿；美国信贷市场规模（2002年）约\$29万亿，其中金融部门借款占\$9万亿；对照之下，美国股票总市值（2002年）约\$12万亿——债券市场的体量明显大于股票市场。

衍生品有两条交易渠道：受监管的交易所与不受监管的场外市场（over-the-counter，OTC）。市场规模通常以合约的名义价值（notional value）度量。交易所交易衍生品（涵盖各大期权与期货交易所的全部合约）的近期估计为名义金额\$13万亿至\$14万亿；OTC衍生品为特定客户定制，2000年底未平仓合约的名义金额估计为\$95.2万亿，且专家认为此数很可能还偏低。

\section{风险与收益}

本节把投资决策的两要素——收益与风险——变成可计算的量：收益用比例变化定义，风险用收益率的（隐含）概率分布的二阶矩定义。

设在时刻$t$以价格$S(t)$买入某股票，持有时长$T$，期末价格为$S(t+T)$，期间获得红利$d$，则该时段的（比例）收益率（fractional rate of return）$R$定义为
\begin{equation*}
R=\frac{S(t+T)+d-S(t)}{S(t)},
\end{equation*}
而$R/T$为瞬时收益率。风险来自未来股价可能上涨也可能下跌这一不确定性：市场可能繁荣（股价上涨）、衰退（股价暴跌）或沉闷（股价微动）。投资者可以为每种情景（scenario）指派概率，从而把市场的不确定性变成可把握的对象。原书表2.1给出一个典型例子（证券现价\$100，年度情景）：

\begin{center}
\begin{tabular}{llccc}
\hline
$s$ & 情景 & 期末价 & 概率$p(s)$ & 年收益$R(s)$\\
\hline
1 & 沉闷（doldrum） & \$100 & 0.70 & 0.00\\
2 & 下行（downturn） & \$85 & 0.20 & $-0.15$\\
3 & 繁荣（boom） & \$105 & 0.10 & 0.05\\
\hline
\end{tabular}
\end{center}

以离散变量$s$标记情景，$p(s)$为其概率，$R(s)$为其收益。投资的期望收益（expected return）即收益的平均（均值）值：
\begin{equation*}
\bar R=\text{期望收益}=\sum_s p(s)\,R(s)\equiv E[R],
\end{equation*}
其中记号$E[X]$表示随机量$X$的期望值。获得期望收益所冒的风险，显然是收益可能偏离该期望值：概率论告诉我们“实际收益$=$期望收益$\pm$标准差（以一定概率）”。偏离的精确幅度与可能性，只有知道股价$S(t)$的概率分布$p(S)$才能算出。标准差$\sigma$是方差的平方根，方差定义为
\begin{equation*}
\sigma^{2}=\sum_s p(s)\bigl(R(s)-\bar R\bigr)^{2}\equiv E\bigl[(R-\bar R)^{2}\bigr],
\end{equation*}
投资的风险即由$\sigma$度量：风险越大$\sigma$越大，反之亦然。按表2.1的数字，年初投资\$100的投资者年末的期望现金为$\$100\times[1+\bar R]\pm6=\$97.50\pm6$。

若收益分布重尾到方差不存在——例如证券服从Lévy分布——则$\sigma$为无穷，此时更合适的风险度量是在险价值（value at risk）。

接着建立无风险基准。银行定期存款等工具被视为无风险（risk-free）的：无风险工具在时刻$t$的收益率，是存入无风险银行的瞬时存款所赚的比率，称为即期利率（spot interest rate）或隔夜拆借利率，记作$r(t)$。于是对无风险工具
\begin{equation*}
\sigma_{\text{risk-free}}=0,\qquad \text{无风险收益率}=\text{即期利率}=r.
\end{equation*}
风险中性（risk-neutral）投资者满足于等于即期利率$r$的回报。但对风险投资$\sigma>0$，要诱导投资者承担高风险，必须有相应的高回报：资本市场为高风险投资标出溢价，以促成资本流向高风险处。风险溢价（risk premium）是高风险投资的收益率高出无风险利率的数额；对平均年收益率为$\bar R$的投资，原书把比值$(\bar R-r)/\sigma$称为风险溢价，又名Sharpe ratio。\footnote{通行术语与此略有差别：risk premium通常专指分子$\bar R-r$，而比值$(\bar R-r)/\sigma$才是Sharpe ratio。原书二者混用，导读照录并提醒读者注意。}投机者（speculator）只有在分析显示潜在回报带有足够风险溢价时才投入高风险证券——这不同于没有溢价也承担高风险的赌徒（gambler）。

本节末尾，无套利原理（principle of no arbitrage）首次亮相：任何无风险金融工具的收益率都不可能高于无风险利率。换言之，天下没有免费午餐——想收获高回报，就必须承担相应的高风险。其数学含义留待2.5节展开。

\section{货币的时间价值}

本节的任务是回答：不同时刻的美元不能直接相加比较，换算的“汇率”是什么？答案是即期利率，而随机利率下的换算规则由本章第一个编号公式(2.1)给出。

先看钱的价值为何依赖时间。资本的货币形式背后是会随时间折蚀的真实生产性资产；通货膨胀、货币贬值等因素也使金融资产所代表的价值随时间变化；金融资产代表启动或促成经济活动的能力，而机会总与变动不居的环境相连。诸如此类的原因使货币的有效价值强烈依赖时间。

如何估算货币的时间价值？原书的经济学答案是：对货币时间价值的全部内生与外生影响，都汇总在一个量里——即期利率。投入其他风险工具的钱估值更复杂，因为风险溢价因投资者而异；归根结底，货币的时间估值就是把未来值贴现（discounting）成“期望的”现值（present value）。

具体机制：设$t_0$时刻持有\$1，将其投入固定存款之类的无风险工具，并把收益不断再投入同一存款（复利（compounding））。由于无风险工具的收益率为$r(t)$，到未来时刻$t_*$这笔投资的无风险价值变为
\begin{equation*}
t_0<t_*:\quad \$1\ \text{在}\ t_0\ \text{时刻}
=\$\,e^{\int_{t_0}^{t_*}dt\,r(t)}\ \text{在}\ t_*\ \text{时刻}\quad(\text{现金的增值})，
\end{equation*}
反过来，把未来值贴现回今天：
\begin{equation*}
\$1\ \text{在}\ t_*\ \text{时刻}
=\$\,e^{-\int_{t_0}^{t_*}dt\,r(t)}\ \text{在}\ t_0\ \text{时刻}\quad(\text{现金的贴现值})。
\end{equation*}
货币时间价值的实质因此是：\$1并不是衡量价值的正确单位——它的有效价值随时间不断变动；正确的单位是贴现量（discounted quantity）。未来现金流应乘以贴现因子（discount factor）$\exp\{-\int_{t_0}^{t_*}dt\,r(t)\}$折到今天，今天的现金流应乘以它的倒数推到未来。

$r(t)$本身如何决定？从第一性原理出发须研究宏观基本面、货币的供给与需求等；利率反映消费的边际效用——诱使人们放弃当下消费、把钱储蓄（投资）起来备作未来消费的比率。后文研究即期利率模型时，$r(t)$将被当作随机（random）变量处理；此时贴现要对随机函数$r(t)$的一切可能取值取平均，贴现因子成为
\begin{equation*}
t_0<t_*:\quad \$1\ \text{在}\ t_*\ \text{时刻}
=\$\,E\!\left[e^{-\int_{t_0}^{t_*}dt\,r(t)}\right]\ \text{在}\ t_0\ \text{时刻}
\tag{2.1}
\end{equation*}
（即现金的贴现值；式中前置的\$表示金额以美元计，$E$为期望。）

推导主线只有一步，却极为要紧：确定性情形的贴现式$e^{-\int r}$——承认$r(t)$是随机过程——“对贴现因子整体取期望”得到式(2.1)。注意期望作用在整个指数上，而不是“先对$r$取平均再代入指数”。这一差别（Jensen不等式的效应）意味着随机利率下的现值一般不等于用平均利率算出的现值；更重要的是，式(2.1)已经埋下伏笔：期望总要对着某个概率测度取，而“选哪个测度”正是2.5节的主题。式(2.1)也是本章后文零息国债定价公式(2.6)的种子。

\section{无套利、鞅与风险中性测度}

本节是全章的理论核心，要解决的问题是：给衍生品定价时，怎样保证定出的价格本身不制造套利机会？答案可以压缩成一句话：在某个（风险中性）测度下让贴现价格成为鞅，然后“价格$=$贴现后的条件期望”。

套利（arbitrage）的定义：通过同时进入两笔或多笔金融交易——无论在同一市场还是跨市场——获取无风险（有保证）的利润。既然经济中存在现金存款之类的无风险工具，套利就特指获得高于货币市场无风险收益的保证收益。原书举了一个教科书式例子：同一公司股票在纽约交易所值US\$1，在新加坡值S\$1.8，而汇率为US\$1$=$S\$1.7；经纪人同时在纽约买进100股、在新加坡卖出100股，即可锁定S\$10的无风险利润。交易成本会消灭小散户的套利机会，但对交易成本近乎为零的大型券商，套利是重要的利润来源；同时套利者的抛售又把新加坡股价压向S\$1.7——套利活动本身就是价格回归均衡的动力。

有效市场与无套利的关系要说清楚：有效市场中不存在套利机会；套利是资本市场在实践中得以作为有效市场运转、并确定金融工具均衡（“正确”）价格的机制之一。有效市场的存在是无套利原理成立的充分条件而非必要条件；均衡中无套利机会成立。无套利是一个稳健（robust）的概念，因为它只表达“所有投资者宁要多财富、不要少财富”这一偏好；而大多数市场均衡模型对投资者行为施加了更苛刻的假设。

理论金融的一个重要结果（原书引[23,40,100]）是：要使金融工具的价格杜绝任何套利可能，就必须让该工具的贴现价值按鞅过程（martingale process）演化。要注意：证券（如股票）的真实市场演化并不是鞅过程——若是，持有该证券便没有风险溢价。鞅演化是给衍生品定价而设的理论构造：先在“风险中性世界”里定价，使价格免于套利。概率论中的鞅（原书附录A.1讨论）是“公平博弈”（fair game）概念的数学化；有效市场中证券的无风险演化等价于其演化满足鞅条件（martingale condition）。由于真实投资者并非风险中性、要求风险溢价，他们的（真实）演化需要从风险中性测度出发作变换测度（change of measure）。

鞅的定义。设随机过程由$N+1$个随机变量$X_i$（$1\le i\le N+1$）给出，联合概率分布函数为$p(x_1,x_2,\ldots,x_{N+1})$（原书附录式(A.2)），则鞅过程由如下条件概率定义：
\begin{equation*}
E\!\left[X_{n+1}\mid x_1,x_2,\ldots,x_n\right]=x_n
\qquad:\ \text{鞅}
\tag{2.2}
\end{equation*}
左端是随机变量$X_{n+1}$在$X_1,X_2,\ldots,X_n$取值为$x_1,x_2,\ldots,x_n$条件下的条件期望（conditional expectation）。在金融应用中，这些随机变量取某股票在时刻$t_1,t_2,\ldots,t_{N+1}$的未来价格$S_1,S_2,\ldots,S_{N+1}$。

要把鞅条件用于股价演化，必须先贴现：被比较的是两个不同时刻的价格，按2.4节的教训须先把金额换到同一时刻的时间可比单位——贴现由无风险利率完成。原书附录A.6证明：对完备市场（complete market），存在唯一（unique）的风险中性测度（risk-neutral measure），使得该资产的一切衍生品的贴现演化都满足鞅条件（引[40,42]）。\footnote{现代教科书的“资产定价基本定理”表述为两条：无套利$\Leftrightarrow$存在等价鞅测度（第一定理）；市场完备$\Leftrightarrow$测度唯一（第二定理）。原书的表述与之一致，参见Harrison--Kreps与Harrison--Pliska的经典谱系及原书附录A.6；对照阅读可参考Shreve《Stochastic Calculus for Finance》。}

定价公式族的出发点：设某股权在未来时刻$t$的价值为$S(t)$，并假设存在贴现股价的无风险演化
\begin{equation*}
e^{-\int_0^t r(t')\,dt'}\,S(t)
\tag{2.3}
\end{equation*}
使它按鞅过程演化（引[42]）。由鞅的条件期望表述（原书附录式(A.40)），股权贴现值在时刻$t$的条件期望就是其现值$S(0)$，换言之
\begin{equation*}
S(0)=E\!\left[\left.e^{-\int_0^t r(t')\,dt'}\,S(t)\,\right|S(0)\right].
\tag{2.4}
\end{equation*}

推导主线值得逐步拆开（这一步没有任何高深数学，力量全部来自测度的选择）：记$Y_t=e^{-\int_0^t r(t')dt'}S(t)$，即把$t$时刻的股价用无风险利率贴现回零时刻；(2.2)中取$n=0$，“以今日全部信息为条件，$Y_t$的期望等于$Y_0$”，而$Y_0=S(0)$，这正是式(2.4)。换个角度读：式(2.4)说“贴现后的未来价格的期望$=$今天的价格”。为什么这等价于无套利？因为若期望价格偏离$S(0)$，买低卖高并持有到期即可获得超出无风险利率的保证收益——恰是套利的定义；鞅条件堵死了这条路。换测度不改变路径的实现，只改变路径的概率权重，从而把真实世界漂移（含风险溢价）换成无风险漂移，这就是“风险中性定价”的机制。

式(2.4)的通用性极强——它对任何证券都成立，是全书引用频率最高的公式之一：给衍生品定价时，把$S(t)$换成衍生品价格、在风险中性测度下取贴现期望即可；本章的式(2.6)与(2.10)都是它的直接特例。有作者（引[42]）把鞅测度的存在性称为金融学基本定理（fundamental theorem of finance），原书附录A.6有简述。总结成一对双向命题：若存在测度使贴现工具的演化满足鞅条件，则该工具一切衍生品的价格免于套利机会；反之，若贴现价格免于套利机会，则存在鞅测度。本书所分析的大多数模型都在鞅测度下演化，从而保证一切（衍生）金融工具的价格无套利。

\section{对冲}

本节处理一个自然的问题：金融工具的演化既然随机，能否用若干风险资产拼出一个无风险的投资组合（portfolio）？换言之，能否用一个工具的随机涨落抵消另一个工具的随机涨落，且抵消精确到组合整体成为无风险证券？对冲还有一层比较功能：收益相同的两个组合，经过对冲的那个风险更低，因而一般是更优的组合。

对冲（hedging）是把某金融工具放入与其他相关工具构成的投资组合、以削减其随机涨落的程序的统称。完全对冲（perfectly hedged）的组合免除一切随机涨落——被对冲工具价格的随机涨落恰被组合内其他工具的反向涨落精确抵消；实践中通常最好的结果也只是部分对冲（partially hedged）。

对冲类似买保险：其成本是买卖所需证券的交易成本，是降险必须支付的价格；交易成本高使对冲更贵，但它对抗风险依然有效。对冲常有的副作用是压低未来收益，原书举国债对冲为例：期货合约入场零成本，可用期货空头对冲债券——利率上升时对冲起效（期货盈利抵补债券跌价）；但利率下降时债券涨价，期货却亏损，反而拉低净利润。消灭“坏”波动的同时也消灭了“好”波动。期权入场虽非零成本，却常能让投资者管理风险而不必接受未来收益缩水。总之，对冲策略取决于投资者的目标。

并非一切涨落皆可对冲。完备市场中原则上存在可用于对冲某一特定工具每种风险的资产；实践中能否完全对冲取决于市场上实际可得的其他工具——衍生工具发展的主要动力，正来自对常用金融工具的对冲需求。

对冲的结构条件：至少需要第二种工具，才能尝试两种工具涨落之间的抵消；这第二种工具必须依赖被对冲的工具，否则二者的随机涨落谈不上关联。对冲初级工具时常需要其衍生工具，反之亦然：衍生工具与标的由同一随机过程驱动，其演化与标的完全相关（perfectly correlated），从而允许完全对冲。

看具体构造。设某证券（如普通股）由随机变量$S(t)$代表，要削弱持有$S(t)$的风险，需要持有第二工具——$S(t)$的一个衍生品，记作$D(S)$——使组合整体的涨落变小。设$S(t)$可以被完全对冲，对冲组合记作$\Pi(S,t)$；组合的具体形态可以是：做多（买入）一份衍生品$D(S)$，同时卖空价值$\Delta(S)$的标的股票$S$。于是在时刻$t$
\begin{equation*}
\Pi(S,t)=D(S)-\Delta(S)\,S .
\end{equation*}
由于证券$S(t)$在时刻$t$的值是已知的，若组合$\Pi(S,t)$的时间演化不含任何随机涨落——实为确定性函数——则组合被完全对冲：
\begin{equation*}
\frac{d\Pi(S,t)}{dt}:\ \text{无随机性}\ \Longrightarrow\ \text{完全对冲}.
\end{equation*}
$\mathrm{d}\Pi(S,t)/\mathrm{d}t$既无随机涨落，它就是无风险证券；无套利原理随即要求完全对冲组合的收益率等于无风险即期利率$r(t)$：
\begin{equation*}
\frac{d\Pi(S,t)}{dt}=r(t)\,\Pi(S,t).
\end{equation*}

推导主线（这是Black--Scholes--Merton论证的骨架，值得记牢）：选对冲参数$\Delta(S)$使得组合瞬时无风险——无套利强制确定性组合只赚无风险利率——由此反解出衍生品$D(S)$必须满足的方程。第3章将用伊藤微积分展开$\mathrm{d}\Pi/\mathrm{d}t$、定出$\Delta$并导出Black--Scholes方程（含随机波动率时为Merton--Garman方程）。\footnote{$\Delta(S)$就是期权交易员口中的Delta；原书此处把它抽象为“卖空价值$\Delta(S)$的股票”，第3章给出其显式表达式。}

实践中对冲可行还须满足若干条件（原书列了三条）：其一，市场必须交易衍生品$D(S)$，否则无法构造对冲组合——许多金融工具因缺乏适当的衍生工具而无法对冲，例如证券的波动率本身。其二，须查明对冲参数（hedging parameter）$\Delta(S)$是否存在、其对股价$S$的函数依赖为何——这要求知道$D(S)$与$S(t)$的精确关系，以及对$S(t)$随机动力学的完整刻画。其三，组合$\Pi(S,t)$依赖时间，对冲必须连续进行——这要求市场有足够流动性，流动性又决定对冲的交易成本。股权对冲的概念在第3章（衍生品）展开，国债对冲在第9章详细讨论。

\section{远期利率与固定收益证券}

本节把视野从股权（一维随机对象$S(t)$）切换到债务市场：那里的随机对象是整条利率曲线。远期利率与固定收益证券是债务市场的基本构件；原书把它们写成“用远期利率逐段贴现\$1”，由此得到三个恒等式(2.5)、(2.8)、(2.9)与一个鞅定价式(2.10)——这正是后半书把远期利率升级为量子场的出发点。

即期利率再回顾：在时刻$t$的瞬时贷款，借款成本是即期利率$r(t)$，通常按年百分比报价；即期利率的典型范围是每年$0.1\%$至$20\%$。

远期利率（forward interest rate，forward rate）的定义：借款人常要在未来特定时刻借钱——例如为一年后的一笔商品买卖——希望现在就锁定届时利率。资本市场为此提供的工具就是远期利率，记作$f(t,x)$：在较早时刻$t<x$以合约形式约定的、未来时段$x$与$x+dx$之间借款的瞬时利率。从数学观点看，即期利率与远期利率都是借钱——借一个无穷小时间$\epsilon$——的瞬时成本。全体远期利率构成利率的期限结构（term structure），并与收益率曲线（yield curve）相关。由定义，时刻$t$的隔夜贷款的即期利率为
\begin{equation*}
r(t)=f(t,t).
\tag{2.5}
\end{equation*}

债券。债券是政府与公司为从资本市场融资而发行的债务工具，发行人承担支付一列预定固定现金流的义务，故各类债券通称固定收益证券（fixed income securities）。国债（Treasury Bond）是无违约风险的工具；公司债、市政债及某些国家（如俄罗斯、阿根廷）的主权债原则上存在违约风险。美债的无风险性质，使美国政府能以尽可能低的利率进行大规模国际举债。零息国债（zero coupon Treasury Bond）是无风险金融工具，只有单一现金流：在未来时刻$T$支付\$1的固定收益；其在$t<T$时刻的价格记作$P(t,T)$，且$P(T,T)=1$。

定价路线一：用即期利率。由货币的时间价值，到期时刻为$T$的债券在到期前的价值，是把$P(T,T)=1$用即期利率贴现到$t$；利率被视为随机变量的一般情形下，按式(2.1)取期望得
\begin{align*}
P(t,T)&=E\!\left[e^{-\int_t^T dt'\,r(t')}\,P(T,T)\right]\\
&=E\!\left[e^{-\int_t^T dt'\,r(t')}\right]
\tag{2.6}
\end{align*}
（即$P(T,T)=1$的贴现值），其中期望值对$r(t)$所服从的随机过程取。式(2.6)表明：国债价格只是即期利率初值$r(t)$的泛函。

定价路线一（续）：附息债券化归零息债券。附息国债$\mathcal{B}(t,T)$有一列预定现金流：在递增时刻$T_i$支付价值$c_i$的息票（coupon），到期时刻$T$偿还价值$L$的本金（principal）。用“现金流相同的两种金融工具是同一的”这一原则，可以证明（引[58]）$\mathcal{B}(t,T)$可用零息债券表出：
\begin{equation*}
\mathcal{B}(t,T)=\sum_{i=1}^{K}c_i\,P(t,T_i)+L\,P(t,T).
\tag{2.7}
\end{equation*}
可见附息债券等价于一个零息债券的组合：任何零息债券模型自动给出附息债券模型。市政债、公司债与高收益债因税收规则、流动性等因素而更难建模，且有有限的违约概率，带有国债所没有的风险成分，因此须支付高于国债的风险溢价。

定价路线二：用远期利率（推导重点）。零息债券是发行人借入的一笔有限期贷款，而远期利率只对瞬时未来借款有定义，故须用远期利率逐段贴现、迭代得到现值。到期时$P(T,T)=\$1$；把\$1从未来时刻$T$逐段贴现到当前时刻$t$。做法：把时间切成间隔$\epsilon$的格点，远期利率$f(t,x_n)$定义在未来时刻$x_n=t+n\epsilon$（$n=0,1,\ldots,[(T-t)/\epsilon]$）。从未来时刻$x_n$到$x_{n-\epsilon}$的瞬时贷款的贴现因子是$e^{-\epsilon f(t,x_n)}$。对时刻$T$的确定性收益\$1逐段贴现到$t$，得
\begin{equation*}
P(t,T)=e^{-\epsilon f(t,x_0)}\,e^{-\epsilon f(t,x_1)}\cdots e^{-\epsilon f(t,x_{N-1})}\,e^{-\epsilon f(t,x_N)}\cdot 1 .
\end{equation*}
取$\epsilon\to 0$的极限：
\begin{align*}
P(t,T)&=e^{-\int_t^T dx\,f(t,x)}
\tag{2.8}\\[2pt]
\Rightarrow\ -\frac{\partial}{\partial T}\ln P(t,T)&=f(t,T)
\tag{2.9}
\end{align*}
式(2.8)用远期利率表出国债价格，是恒等式，可反过来取作远期利率的定义；由式(2.8)与(2.9)，$P(t,T)$已知则$f(t,T)$亦已知，反之亦然。

鞅条件用于远期利率。把式(2.4)的精神用于零息债券的演化：设到期时刻$T$的零息国债在$t_0<T$的价格为$P(t_0,T)$，在$t_0<t_*<T$的价格为$P(t_*,T)$。满足鞅条件意味着：$t_*$时刻的债券价格沿时间倒推至$t_0$、并以无风险即期利率$r(t)$连续贴现后，其（条件）期望必须等于$t_0$时刻的债券价格。由式(2.4)，零息债券的鞅条件有如下非常一般、与模型无关（model-independent）的表达：
\begin{equation*}
P(t_0,T)=E_{[t_0,t_*]}\!\left[\left.e^{-\int_{t_0}^{t_*}r(t)\,dt}\,P(t_*,T)\,\right|P(t_0,T)\right],
\tag{2.10}
\end{equation*}
其中记号$E_{[t_0,t_*]}[X]$表示随机变量$X$在时段$(t_0,t_*]$上的平均。

两种定义的对峙是本节的收束，也是后半书的路线选择。零息债券有两个定义：用即期利率的式(2.6)与用远期利率的式(2.8)，二者很不一样。可以取即期利率为基本量（远期利率为导出量），如式(2.6)：优点是简单，适于处理只与即期利率行为直接相关的问题；缺点是即期利率必须与观察到的各期限债券价格全部相容。而式(2.8)把$P(t,T)$写成远期利率，是更一般的表达式：远期利率原则上可直接由市场确定；对全部期限的远期利率建模时，以整条初始期限结构作为输入，因此模型不会基于观测到的债券价格生成套利机会。代价是$f(t,x)$依赖两个变量$t,x$，而$r(t)$只依赖时间$t$。尽管更复杂，把远期利率视为基本量、把即期利率仅仅当作远期利率曲线上的一点，仍是远期利率建模的一条高产路线。\footnote{这正是HJM模型（Heath--Jarrow--Morton框架）的哲学：直接为整条瞬时远期利率曲线指定随机动力学，并以无套利约束漂移。本书后半书更进一步，把$f(t,x)$写成由量子场$\mathcal{A}(t,x)$驱动的场论模型。}

\section{小结}

原书以几段短论收束全章，每段都是全书路线图的一条支线，导读整理如下。第一，金融工具的随机演化是用随机变量为金融建模的全部根基；正是金融工具的随机时间演化，提供了与量子理论及路径积分的关键连接（第4章起全面展开）。第二，风险概念居于对冲概念的中心；把风险概念推广到无穷维情形，将在为国债对冲建立量子场论（第9章）时起关键作用。第三，货币的时间价值自然引出鞅测度——其定义式的要求是“贴现后的未来值等于现值”（即式(2.4)的形状）；为非线性远期利率找到恰当的鞅测度，将耗费全书大量笔墨。第四，对冲与期权构成金融衍生品理论的基石，期权定价是本书前半的焦点。第五，本书后半处理远期利率与固定收益证券；对远期利率的研究自然引向把它表述为量子场——这是全书的主旨。

读法建议（导读）：把本章10个编号公式分成三组记忆——定价公式族(2.1)、(2.4)、(2.6)、(2.10)同属“贴现期望”结构，式(2.4)是母式，其余是特例；(2.5)、(2.8)、(2.9)是$r$、$f$、$P$三者间的恒等关系，构成利率端的概念字典；(2.2)是唯一的概率论定义，(2.3)是其作用对象，(2.7)则是化归（附息$\to$零息）。带着这张公式地图进入第3章，Black--Scholes方程的每一个组成部分都已在本章就位。

%TAGS: 2.1,2.2,2.3,2.4,2.5,2.6,2.7,2.8,2.9,2.10
